\documentclass[11pt,a4paper]{article}
\usepackage[margin=25mm]{geometry}
\usepackage[T1]{fontenc}
\usepackage{newtxtext,newtxmath}
\usepackage{parskip}
\setlength{\emergencystretch}{2em}
\pagestyle{empty}

\begin{document}
{\Large\bfseries Submission note for \textit{PeerJ Computer Science}\par}
\medskip
\textit{Author-review draft for the submission form.}

\medskip
\textbf{Article type.} Research Article

\textbf{Manuscript.} Three Protocol Choices Shape Measured Fidelity of AI Inpainting on Museum Paintings

\textbf{Authors.} Wei Meng, Wei Xiong, and Zhanhao Liang.

Dear Editors,

Please consider our manuscript for publication in \textit{PeerJ Computer Science}. We examine how three implementation and evaluation choices affect reported results for off-the-shelf image inpainting on digitized paintings: mask downsampling to the latent grid, hole prefilling when denoising begins from an encoded image, and the denominator used to account for outputs withheld by a safety checker.

We compare four fixed models on 120 open-access paintings from the Cleveland Museum of Art and the Art Institute of Chicago, spanning four medium-and-support groups and six procedural damage conditions. These choices can materially change measured error and output availability. Across the two cohorts and all tested media and damage conditions, the controlled model ranking remained stable. The study contributes a reproducible benchmark and practical reporting guidance for researchers who evaluate image restoration systems on digitized cultural collections.

All authors approved the manuscript and agree to its submission to \textit{PeerJ Computer Science}.

\textbf{Author checks before submission.} Confirm the manuscript's publication and review status, disclose any related preprint or earlier version, and verify that reviewers can access the supporting data and code. Confirm each author's drafting or substantive revision role. For the AI-use declaration, record the Gemini and Codex model versions and identify the author responsible for verifying the GPT-6 translation.

Sincerely,

Wei Xiong, corresponding author\\
Guilin Institute of Information Technology\\
xiongwei@guit.edu.cn\\
\textit{[Submission date]}
\end{document}
